\subsection{$L^p$ 上的 Fourier 变换, $1<p\le2$}
\label{sec:ch1-fourier-transform-lp}

若 $f\in L^p$, $1\le p\le\infty$, 则可将 $f$ 等同于缓增分布: 对 $\phi\in\mathcal S$, 定义
\[
 T_f(\phi)=\int_{\mathbb R^n}f\phi.
\]
显然这个积分有限. 为说明 $T_f$ 连续, 设当 $k\to\infty$ 时 $\phi_k\to0$ 于 $\mathcal S$. 由 Hölder 不等式,
\[
 |T_f(\phi_k)|\le\|f\|_p\|\phi_k\|_{p'}.
\]
$\|\phi_k\|_{p'}$ 可由形如 $x^\alpha\phi_k$ 的函数的 $L^\infty$ 范数控制, 从而可由 $\phi_k$ 的有限个半范数的线性组合控制; 因而左端在 $k\to\infty$ 时趋于 $0$.

此外, 当 $1\le p\le2$ 时, $\widehat f$ 是一个函数.

\begin{Theorem}
\label{thm:ch1-plancherel}
Fourier 变换是 $L^2$ 上的等距映射; 即 $\widehat f\in L^2$ 且 $\|\widehat f\|_2=\|f\|_2$. 此外,
\[
 \widehat f(\xi)=\lim_{R\to\infty}\int_{|x|<R}f(x)e^{-2\pi ix\cdot\xi}\,dx,
\]
并且
\[
 f(x)=\lim_{R\to\infty}\int_{|\xi|<R}\widehat f(\xi)e^{2\pi ix\cdot\xi}\,d\xi,
\]
其中极限均在 $L^2$ 中取.
\end{Theorem}

恒等式 $\|\widehat f\|_2=\|f\|_2$ 称为 Plancherel 定理.

\hypertarget{chapter-01-page-028}{}
\begin{Proof}
给定 $f,h\in\mathcal S$, 令 $g=\widetilde{\overline h}$, 从而 $\widehat g=\overline{\widehat h}$. 由 \eqref{eq:ch1-fourier-pairing-symmetry},
\begin{equation}
\label{eq:ch1-plancherel-pairing}
 \int_{\mathbb R^n}f\overline h=\int_{\mathbb R^n}\widehat f\,\overline{\widehat h}.
\end{equation}
令 $h=f$, 得到对 $f\in\mathcal S$ 有 $\|f\|_2=\|\widehat f\|_2$. 因为 $\mathcal S$ 在 $L^2$ 中稠密, Fourier 变换以范数相等的方式延拓到所有 $f\in L^2$.

最后, Fourier 变换的连续性蕴含所述的 $f$ 和 $\widehat f$ 的极限公式, 因为 $f\chi_{B(0,R)}$ 和 $\widehat f\chi_{B(0,R)}$ 分别在 $L^2$ 中收敛到 $f$ 和 $\widehat f$.
\end{Proof}

若 $f\in L^p$, $1<p<2$, 则可分解为 $f=f_1+f_2$, 其中 $f_1\in L^1$, $f_2\in L^2$. 例如可取 $f_1=f\chi_{\{x:|f(x)|>1\}}$, $f_2=f-f_1$. 因此 $\widehat f=\widehat f_1+\widehat f_2\in L^\infty+L^2$. 然而, 应用插值定理可以得到更精确的结果.

\begin{Theorem}[Riesz--Thorin 插值]
\label{thm:ch1-riesz-thorin}
设 $1\le p_0,p_1,q_0,q_1\le\infty$, 并且对 $0<\theta<1$, 由
\[
 \frac1p=\frac{1-\theta}{p_0}+\frac\theta{p_1},
 \qquad
 \frac1q=\frac{1-\theta}{q_0}+\frac\theta{q_1}
\]
定义 $p$ 和 $q$. 若 $T$ 是从 $L^{p_0}+L^{p_1}$ 到 $L^{q_0}+L^{q_1}$ 的线性算子, 且
\[
 \|Tf\|_{q_0}\le M_0\|f\|_{p_0}\quad\text{for }f\in L^{p_0},
\]
以及
\[
 \|Tf\|_{q_1}\le M_1\|f\|_{p_1}\quad\text{for }f\in L^{p_1},
\]
则
\[
 \|Tf\|_q\le M_0^{1-\theta}M_1^\theta\|f\|_p
 \quad\text{for }f\in L^p.
\]
\end{Theorem}

这一结果的证明使用解析函数的所谓“三线定理”; 例如可见 Stein 与 Weiss \MBUnresolvedReference{citation}{[18, Chapter 5]}, 或 Katznelson \MBUnresolvedReference{citation}{[10, Chapter 4]}.

\begin{Corollary}[Hausdorff--Young 不等式]
\label{cor:ch1-hausdorff-young}
若 $f\in L^p$, $1\le p\le2$, 则 $\widehat f\in L^{p'}$ 且
\[
 \|\widehat f\|_{p'}\le\|f\|_p.
\]
\end{Corollary}

\begin{Proof}
在定理~\ref{thm:ch1-riesz-thorin} 中使用不等式 \eqref{eq:ch1-fourier-l1-linfty}, 即 $\|\widehat f\|_\infty\le\|f\|_1$, 以及 Plancherel 定理 $\|\widehat f\|_2=\|f\|_2$.
\end{Proof}

\hypertarget{chapter-01-page-029}{}
暂且离开 Fourier 变换, 给出 Riesz--Thorin 插值的另一个推论; 它将在后面的章节中使用.

\begin{Corollary}[Young 不等式]
\label{cor:ch1-young-convolution}
若 $f\in L^p$ 且 $g\in L^q$, 则 $f*g\in L^r$, 其中 $1/r+1=1/p+1/q$, 并且
\[
 \|f*g\|_r\le\|f\|_p\|g\|_q.
\]
\end{Corollary}

\begin{Proof}
固定 $f\in L^p$, 立即得到不等式
\[
 \|f*g\|_p\le\|f\|_p\|g\|_1
\]
以及
\[
 \|f*g\|_\infty\le\|f\|_p\|g\|_{p'}.
\]
所需结果由 Riesz--Thorin 插值得到.
\end{Proof}

